\documentclass[12pt,a4paper]{article}
\usepackage[utf8]{inputenc}
\usepackage[T1]{fontenc}
\usepackage[indonesian]{babel}
\usepackage[a4paper,margin=2.5cm]{geometry}
\usepackage{array}
\usepackage{booktabs}
\usepackage{helvet}
\usepackage[hidelinks]{hyperref}
\usepackage{xcolor}

\renewcommand{\familydefault}{\sfdefault}
\setlength{\parindent}{0pt}
\setlength{\parskip}{0.7em}

\begin{document}

\begin{center}
{\LARGE \textbf{Informasi Kuis Sistem Operasi}}\\[0.4em]
{\large Migrasi Akses Kuis ke Moodle}
\end{center}

Mahasiswa dimohon menggunakan Moodle untuk akses kuis mata kuliah ini. Migrasi ini dilakukan sebagai langkah antisipasi apabila \textit{lectura.id} mengalami \textit{maintenance}, mengingat kredit di sisi server saat ini terbatas.

\textbf{Link course Moodle:}\\
\href{https://kuliah2.itera.ac.id/course/view.php?id=656}{https://kuliah2.itera.ac.id/course/view.php?id=656}

\textbf{Enrollment key per kelas:}

\begin{center}
\begin{tabular}{>{\bfseries}c l}
\toprule
Kelas & Enrollment key \\
\midrule
RA & \texttt{@OSra2026} \\
RB & \texttt{@OSrb2026} \\
RC & \texttt{@OSrc2026} \\
RD & \texttt{@OSrd2026} \\
RE & \texttt{@OSre2026} \\
RF & \texttt{@OSrf2026} \\
\bottomrule
\end{tabular}
\end{center}

\textbf{Catatan:}
\begin{itemize}
    \item Pastikan masuk ke course Moodle yang benar sebelum mengerjakan kuis.
    \item Simpan enrollment key sesuai kelas masing-masing dan jangan tertukar.
    \item Jika sebelumnya mengakses kuis melalui \textit{lectura.id}, silakan beralih ke Moodle untuk menghindari kendala saat sistem sedang \textit{maintenance}.
\end{itemize}

Jika ada kendala akses atau masalah teknis lainnya, silakan langsung chat:\\
\href{https://wa.me/62881080130131}{wa.me/62881080130131} (Pak MCT)

\vfill
\begin{flushright}
Sistem Operasi\\
Semester Genap 2025/2026
\end{flushright}

\end{document}